% Options for packages loaded elsewhere
% Options for packages loaded elsewhere
\PassOptionsToPackage{unicode}{hyperref}
\PassOptionsToPackage{hyphens}{url}
\PassOptionsToPackage{dvipsnames,svgnames,x11names}{xcolor}
%
\documentclass[
  letterpaper,
  DIV=11,
  numbers=noendperiod]{scrartcl}
\usepackage{xcolor}
\usepackage[margin=1in]{geometry}
\usepackage{amsmath,amssymb}
\setcounter{secnumdepth}{5}
\usepackage{iftex}
\ifPDFTeX
  \usepackage[T1]{fontenc}
  \usepackage[utf8]{inputenc}
  \usepackage{textcomp} % provide euro and other symbols
\else % if luatex or xetex
  \usepackage{unicode-math} % this also loads fontspec
  \defaultfontfeatures{Scale=MatchLowercase}
  \defaultfontfeatures[\rmfamily]{Ligatures=TeX,Scale=1}
\fi
\usepackage{lmodern}
\ifPDFTeX\else
  % xetex/luatex font selection
  \setmonofont[]{DejaVu Sans Mono}
\fi
% Use upquote if available, for straight quotes in verbatim environments
\IfFileExists{upquote.sty}{\usepackage{upquote}}{}
\IfFileExists{microtype.sty}{% use microtype if available
  \usepackage[]{microtype}
  \UseMicrotypeSet[protrusion]{basicmath} % disable protrusion for tt fonts
}{}
\makeatletter
\@ifundefined{KOMAClassName}{% if non-KOMA class
  \IfFileExists{parskip.sty}{%
    \usepackage{parskip}
  }{% else
    \setlength{\parindent}{0pt}
    \setlength{\parskip}{6pt plus 2pt minus 1pt}}
}{% if KOMA class
  \KOMAoptions{parskip=half}}
\makeatother
% Make \paragraph and \subparagraph free-standing
\makeatletter
\ifx\paragraph\undefined\else
  \let\oldparagraph\paragraph
  \renewcommand{\paragraph}{
    \@ifstar
      \xxxParagraphStar
      \xxxParagraphNoStar
  }
  \newcommand{\xxxParagraphStar}[1]{\oldparagraph*{#1}\mbox{}}
  \newcommand{\xxxParagraphNoStar}[1]{\oldparagraph{#1}\mbox{}}
\fi
\ifx\subparagraph\undefined\else
  \let\oldsubparagraph\subparagraph
  \renewcommand{\subparagraph}{
    \@ifstar
      \xxxSubParagraphStar
      \xxxSubParagraphNoStar
  }
  \newcommand{\xxxSubParagraphStar}[1]{\oldsubparagraph*{#1}\mbox{}}
  \newcommand{\xxxSubParagraphNoStar}[1]{\oldsubparagraph{#1}\mbox{}}
\fi
\makeatother

\usepackage{color}
\usepackage{fancyvrb}
\newcommand{\VerbBar}{|}
\newcommand{\VERB}{\Verb[commandchars=\\\{\}]}
\DefineVerbatimEnvironment{Highlighting}{Verbatim}{commandchars=\\\{\}}
% Add ',fontsize=\small' for more characters per line
\usepackage{framed}
\definecolor{shadecolor}{RGB}{241,243,245}
\newenvironment{Shaded}{\begin{snugshade}}{\end{snugshade}}
\newcommand{\AlertTok}[1]{\textcolor[rgb]{0.68,0.00,0.00}{#1}}
\newcommand{\AnnotationTok}[1]{\textcolor[rgb]{0.37,0.37,0.37}{#1}}
\newcommand{\AttributeTok}[1]{\textcolor[rgb]{0.40,0.45,0.13}{#1}}
\newcommand{\BaseNTok}[1]{\textcolor[rgb]{0.68,0.00,0.00}{#1}}
\newcommand{\BuiltInTok}[1]{\textcolor[rgb]{0.00,0.23,0.31}{#1}}
\newcommand{\CharTok}[1]{\textcolor[rgb]{0.13,0.47,0.30}{#1}}
\newcommand{\CommentTok}[1]{\textcolor[rgb]{0.37,0.37,0.37}{#1}}
\newcommand{\CommentVarTok}[1]{\textcolor[rgb]{0.37,0.37,0.37}{\textit{#1}}}
\newcommand{\ConstantTok}[1]{\textcolor[rgb]{0.56,0.35,0.01}{#1}}
\newcommand{\ControlFlowTok}[1]{\textcolor[rgb]{0.00,0.23,0.31}{\textbf{#1}}}
\newcommand{\DataTypeTok}[1]{\textcolor[rgb]{0.68,0.00,0.00}{#1}}
\newcommand{\DecValTok}[1]{\textcolor[rgb]{0.68,0.00,0.00}{#1}}
\newcommand{\DocumentationTok}[1]{\textcolor[rgb]{0.37,0.37,0.37}{\textit{#1}}}
\newcommand{\ErrorTok}[1]{\textcolor[rgb]{0.68,0.00,0.00}{#1}}
\newcommand{\ExtensionTok}[1]{\textcolor[rgb]{0.00,0.23,0.31}{#1}}
\newcommand{\FloatTok}[1]{\textcolor[rgb]{0.68,0.00,0.00}{#1}}
\newcommand{\FunctionTok}[1]{\textcolor[rgb]{0.28,0.35,0.67}{#1}}
\newcommand{\ImportTok}[1]{\textcolor[rgb]{0.00,0.46,0.62}{#1}}
\newcommand{\InformationTok}[1]{\textcolor[rgb]{0.37,0.37,0.37}{#1}}
\newcommand{\KeywordTok}[1]{\textcolor[rgb]{0.00,0.23,0.31}{\textbf{#1}}}
\newcommand{\NormalTok}[1]{\textcolor[rgb]{0.00,0.23,0.31}{#1}}
\newcommand{\OperatorTok}[1]{\textcolor[rgb]{0.37,0.37,0.37}{#1}}
\newcommand{\OtherTok}[1]{\textcolor[rgb]{0.00,0.23,0.31}{#1}}
\newcommand{\PreprocessorTok}[1]{\textcolor[rgb]{0.68,0.00,0.00}{#1}}
\newcommand{\RegionMarkerTok}[1]{\textcolor[rgb]{0.00,0.23,0.31}{#1}}
\newcommand{\SpecialCharTok}[1]{\textcolor[rgb]{0.37,0.37,0.37}{#1}}
\newcommand{\SpecialStringTok}[1]{\textcolor[rgb]{0.13,0.47,0.30}{#1}}
\newcommand{\StringTok}[1]{\textcolor[rgb]{0.13,0.47,0.30}{#1}}
\newcommand{\VariableTok}[1]{\textcolor[rgb]{0.07,0.07,0.07}{#1}}
\newcommand{\VerbatimStringTok}[1]{\textcolor[rgb]{0.13,0.47,0.30}{#1}}
\newcommand{\WarningTok}[1]{\textcolor[rgb]{0.37,0.37,0.37}{\textit{#1}}}

\usepackage{longtable,booktabs,array}
\newcounter{none} % for unnumbered tables
\usepackage{calc} % for calculating minipage widths
% Correct order of tables after \paragraph or \subparagraph
\usepackage{etoolbox}
\makeatletter
\patchcmd\longtable{\par}{\if@noskipsec\mbox{}\fi\par}{}{}
\makeatother
% Allow footnotes in longtable head/foot
\IfFileExists{footnotehyper.sty}{\usepackage{footnotehyper}}{\usepackage{footnote}}
\makesavenoteenv{longtable}
\usepackage{graphicx}
\makeatletter
\newsavebox\pandoc@box
\newcommand*\pandocbounded[1]{% scales image to fit in text height/width
  \sbox\pandoc@box{#1}%
  \Gscale@div\@tempa{\textheight}{\dimexpr\ht\pandoc@box+\dp\pandoc@box\relax}%
  \Gscale@div\@tempb{\linewidth}{\wd\pandoc@box}%
  \ifdim\@tempb\p@<\@tempa\p@\let\@tempa\@tempb\fi% select the smaller of both
  \ifdim\@tempa\p@<\p@\scalebox{\@tempa}{\usebox\pandoc@box}%
  \else\usebox{\pandoc@box}%
  \fi%
}
% Set default figure placement to htbp
\def\fps@figure{htbp}
\makeatother





\setlength{\emergencystretch}{3em} % prevent overfull lines

\providecommand{\tightlist}{%
  \setlength{\itemsep}{0pt}\setlength{\parskip}{0pt}}



 


\KOMAoption{captions}{tableheading}
\makeatletter
\@ifpackageloaded{caption}{}{\usepackage{caption}}
\AtBeginDocument{%
\ifdefined\contentsname
  \renewcommand*\contentsname{Table of contents}
\else
  \newcommand\contentsname{Table of contents}
\fi
\ifdefined\listfigurename
  \renewcommand*\listfigurename{List of Figures}
\else
  \newcommand\listfigurename{List of Figures}
\fi
\ifdefined\listtablename
  \renewcommand*\listtablename{List of Tables}
\else
  \newcommand\listtablename{List of Tables}
\fi
\ifdefined\figurename
  \renewcommand*\figurename{Figure}
\else
  \newcommand\figurename{Figure}
\fi
\ifdefined\tablename
  \renewcommand*\tablename{Table}
\else
  \newcommand\tablename{Table}
\fi
}
\@ifpackageloaded{float}{}{\usepackage{float}}
\floatstyle{ruled}
\@ifundefined{c@chapter}{\newfloat{codelisting}{h}{lop}}{\newfloat{codelisting}{h}{lop}[chapter]}
\floatname{codelisting}{Listing}
\newcommand*\listoflistings{\listof{codelisting}{List of Listings}}
\makeatother
\makeatletter
\makeatother
\makeatletter
\@ifpackageloaded{caption}{}{\usepackage{caption}}
\@ifpackageloaded{subcaption}{}{\usepackage{subcaption}}
\makeatother
\makeatletter
\@ifpackageloaded{tcolorbox}{}{\usepackage[skins,breakable]{tcolorbox}}
\makeatother
\makeatletter
\@ifundefined{shadecolor}{\definecolor{shadecolor}{rgb}{.97, .97, .97}}{}
\makeatother
\makeatletter
\makeatother
\makeatletter
\ifdefined\Shaded\renewenvironment{Shaded}{\begin{tcolorbox}[boxrule=0pt, breakable, enhanced, frame hidden, interior hidden, sharp corners]}{\end{tcolorbox}}\fi
\makeatother
\usepackage{bookmark}
\IfFileExists{xurl.sty}{\usepackage{xurl}}{} % add URL line breaks if available
\urlstyle{same}
\hypersetup{
  pdftitle={Setting Up Git for Data Science Workflows},
  pdfauthor={Ronald `Ryy' G. Thomas},
  colorlinks=true,
  linkcolor={blue},
  filecolor={Maroon},
  citecolor={Blue},
  urlcolor={Blue},
  pdfcreator={LaTeX via pandoc}}


\title{Setting Up Git for Data Science Workflows}
\usepackage{etoolbox}
\makeatletter
\providecommand{\subtitle}[1]{% add subtitle to \maketitle
  \apptocmd{\@title}{\par {\large #1 \par}}{}{}
}
\makeatother
\subtitle{A practical guide to version control for solo analysts and
small research teams}
\author{Ronald `Ryy' G. Thomas}
\date{2026-02-11}
\begin{document}
\maketitle

\renewcommand*\contentsname{Table of contents}
{
\setcounter{tocdepth}{3}
\tableofcontents
}

\begin{figure}[H]

{\centering 
\includegraphics[width=0.8\linewidth,height=\textheight,keepaspectratio]{media/images/hero.png}

}

\caption{Git logo and branching diagram representing version control
workflows}

\end{figure}%

\emph{Version control is the foundation of every reproducible research
workflow.}

\section{Introduction}\label{introduction}

I did not really know how much damage a missing backup could cause until
I accidentally overwrote three days of analysis work on an Alzheimer's
disease dataset. The file was gone, the changes were irreversible, and
the only copy lived on a USB drive I had left at the office. That
experience convinced me that version control was not optional for data
science---it was essential.

Git is the most widely adopted version control system in both software
engineering and data science. It tracks every change to every file,
allows branching and merging of parallel work, and provides a complete
history of a project's evolution. Combined with GitHub as a remote
hosting service, Git enables collaboration, code review, and
reproducibility at a level that manual file management simply cannot
match.

This post documents my own path from ``I'll just save copies'' to a
reliable Git-based workflow. It covers installation, basic commands,
branching, collaboration, and the commit message practices that I found
most valuable. I am writing as a learner, not as an expert, and I
welcome corrections or suggestions from readers who have found better
approaches.

\subsection{Motivations}\label{motivations}

\begin{itemize}
\tightlist
\item
  I lost critical analysis files more than once because I relied on
  manual backups and naming conventions like
  \texttt{analysis\_v3\_FINAL\_FINAL.R}.
\item
  I needed a way to experiment with new modeling approaches without
  risking the working version of my code.
\item
  Collaborating with team members on shared analysis projects required a
  system more robust than emailing files back and forth.
\item
  I wanted a clear, timestamped record of every change I made to a
  project, so that I could trace when and why a result changed.
\item
  I was increasingly working across multiple machines (office
  workstation, laptop, server) and needed a single source of truth for
  every project.
\item
  The biostatistics community has been moving toward reproducible
  research standards, and Git is a core component of that movement
  (Bryan, 2018).
\end{itemize}

\subsection{Objectives}\label{objectives}

\begin{enumerate}
\def\labelenumi{\arabic{enumi}.}
\tightlist
\item
  Install and configure Git on macOS or Linux, including identity setup
  and SSH authentication.
\item
  Learn the core Git workflow: staging, committing, pushing, and pulling
  changes.
\item
  Understand branching and merging for safe experimentation and
  collaboration.
\item
  Adopt commit message best practices drawn from the Angular project and
  the imperative mood convention.
\end{enumerate}

I am documenting my learning process here. If you spot errors or have
better approaches, please let me know in the Let's Connect section at
the end.

\begin{figure}[H]

{\centering 
\includegraphics[width=0.3\linewidth,height=\textheight,keepaspectratio]{media/images/placeholder-coffee-04.jpg}

}

\caption{Settling in for version control setup.}

\end{figure}%

\emph{GitHub provides the remote hosting layer that makes Git-based
collaboration practical for research teams.}

\section{Prerequisites and Setup}\label{prerequisites-and-setup}

This tutorial assumes a macOS or Linux environment. The commands shown
use a standard terminal (Bash or Zsh). No prior Git experience is
required, but familiarity with basic command-line navigation
(\texttt{cd}, \texttt{ls}, \texttt{mkdir}) will be helpful.

\textbf{Required software:}

\begin{itemize}
\tightlist
\item
  Git (version 2.30 or later recommended)
\item
  A terminal emulator (Terminal.app, iTerm2, Kitty, or similar)
\item
  A GitHub account (free tier is sufficient)
\item
  A text editor (Vim, Neovim, VS Code, or RStudio)
\end{itemize}

\textbf{Installing Git on macOS:}

\begin{Shaded}
\begin{Highlighting}[]
\ExtensionTok{xcode{-}select} \AttributeTok{{-}{-}install}
\end{Highlighting}
\end{Shaded}

This installs the Apple Command Line Tools, which include Git. To verify
the installation:

\begin{Shaded}
\begin{Highlighting}[]
\FunctionTok{git} \AttributeTok{{-}{-}version}
\end{Highlighting}
\end{Shaded}

\textbf{Installing Git on Linux (Debian/Ubuntu):}

\begin{Shaded}
\begin{Highlighting}[]
\FunctionTok{sudo}\NormalTok{ apt update }\KeywordTok{\&\&} \FunctionTok{sudo}\NormalTok{ apt install git}
\end{Highlighting}
\end{Shaded}

\section{What is Git?}\label{what-is-git}

Git is a distributed version control system. In plain terms, it is a
tool that records every change you make to a set of files, stores those
changes in a structured history, and allows you to retrieve any previous
version of any file at any time. Think of it as an unlimited ``undo''
system that also tracks who made each change and why.

The ``distributed'' part means that every collaborator has a complete
copy of the project history on their own machine. There is no single
server that must be online for work to continue. GitHub, GitLab, and
Bitbucket serve as convenient central repositories for sharing work, but
they are not strictly required for Git itself to function.

For a solo data scientist, Git provides a safety net: every experiment,
every model iteration, and every data cleaning step is preserved. For a
research team, it provides a coordination mechanism that eliminates the
chaos of emailed files and shared drives.

\section{Getting Started: Initial
Configuration}\label{getting-started-initial-configuration}

Before using Git for the first time, you need to set your identity. Git
attaches your name and email to every commit you make.

\begin{Shaded}
\begin{Highlighting}[]
\FunctionTok{git}\NormalTok{ config }\AttributeTok{{-}{-}global}\NormalTok{ user.name }\StringTok{"Your Name"}
\FunctionTok{git}\NormalTok{ config }\AttributeTok{{-}{-}global}\NormalTok{ user.email }\StringTok{"you@example.com"}
\end{Highlighting}
\end{Shaded}

I also recommend setting a default branch name and a preferred editor:

\begin{Shaded}
\begin{Highlighting}[]
\FunctionTok{git}\NormalTok{ config }\AttributeTok{{-}{-}global}\NormalTok{ init.defaultBranch main}
\FunctionTok{git}\NormalTok{ config }\AttributeTok{{-}{-}global}\NormalTok{ core.editor }\StringTok{"vim"}
\end{Highlighting}
\end{Shaded}

To verify your configuration:

\begin{Shaded}
\begin{Highlighting}[]
\FunctionTok{git}\NormalTok{ config }\AttributeTok{{-}{-}list}
\end{Highlighting}
\end{Shaded}

\textbf{Setting up SSH authentication with GitHub:}

SSH keys allow you to push and pull from GitHub without entering your
password each time. Generate a key pair:

\begin{Shaded}
\begin{Highlighting}[]
\FunctionTok{ssh{-}keygen} \AttributeTok{{-}t}\NormalTok{ ed25519 }\DataTypeTok{\textbackslash{}}
  \AttributeTok{{-}C} \StringTok{"you@example.com"}
\end{Highlighting}
\end{Shaded}

Then add the public key to your GitHub account under Settings
\textgreater{} SSH and GPG keys. Test the connection:

\begin{Shaded}
\begin{Highlighting}[]
\FunctionTok{ssh} \AttributeTok{{-}T}\NormalTok{ git@github.com}
\end{Highlighting}
\end{Shaded}

A successful response will say ``Hi username! You've successfully
authenticated.''

\section{Creating Your First
Repository}\label{creating-your-first-repository}

There are two common starting points: creating a new repository from
scratch, or cloning an existing one from GitHub.

\textbf{Creating a new local repository:}

\begin{Shaded}
\begin{Highlighting}[]
\FunctionTok{mkdir}\NormalTok{ my{-}analysis}
\BuiltInTok{cd}\NormalTok{ my{-}analysis}
\FunctionTok{git}\NormalTok{ init}
\end{Highlighting}
\end{Shaded}

This creates a hidden \texttt{.git} directory that stores the entire
version history. The directory itself is now a Git repository.

\textbf{Cloning an existing repository:}

\begin{Shaded}
\begin{Highlighting}[]
\FunctionTok{git}\NormalTok{ clone git@github.com:username/repo.git}
\end{Highlighting}
\end{Shaded}

This downloads the full project, including all history, and sets up the
remote connection automatically.

I found that initializing a repository on GitHub first (with a README
and \texttt{.gitignore} for R) and then cloning it locally was the most
reliable approach for new projects.

\section{The Core Workflow: Add, Commit,
Push}\label{the-core-workflow-add-commit-push}

The daily Git workflow involves three steps: staging changes, committing
them to history, and pushing them to a remote repository.

\textbf{Step 1: Check the current state}

\begin{Shaded}
\begin{Highlighting}[]
\FunctionTok{git}\NormalTok{ status}
\end{Highlighting}
\end{Shaded}

This shows which files have been modified, which are staged for commit,
and which are untracked.

\textbf{Step 2: Stage files for commit}

\begin{Shaded}
\begin{Highlighting}[]
\FunctionTok{git}\NormalTok{ add analysis.R}
\FunctionTok{git}\NormalTok{ add data/clean\_data.csv}
\end{Highlighting}
\end{Shaded}

Or, to stage all changed files:

\begin{Shaded}
\begin{Highlighting}[]
\FunctionTok{git}\NormalTok{ add .}
\end{Highlighting}
\end{Shaded}

I prefer staging files individually rather than using
\texttt{git\ add\ .} because it forces me to review what I am
committing. This prevents accidental inclusion of large data files or
sensitive credentials.

\textbf{Step 3: Commit with a message}

\begin{Shaded}
\begin{Highlighting}[]
\FunctionTok{git}\NormalTok{ commit }\AttributeTok{{-}m} \StringTok{"Add linear regression analysis"}
\end{Highlighting}
\end{Shaded}

Each commit is a snapshot of the project at a specific point in time.
The message should explain what changed and why.

\textbf{Step 4: Push to the remote}

\begin{Shaded}
\begin{Highlighting}[]
\FunctionTok{git}\NormalTok{ push origin main}
\end{Highlighting}
\end{Shaded}

This sends your local commits to GitHub, making them available to
collaborators and serving as an off-site backup.

\textbf{Pulling changes from the remote:}

\begin{Shaded}
\begin{Highlighting}[]
\FunctionTok{git}\NormalTok{ pull origin main}
\end{Highlighting}
\end{Shaded}

Always pull before starting new work to ensure you have the latest
version.

\begin{figure}[H]

{\centering \pandocbounded{\includegraphics[keepaspectratio]{media/images/ambiance2.png}}

}

\caption{Soft atmospheric image suggesting focused work}

\end{figure}%

\emph{Establishing a consistent workflow early prevents most version
control headaches downstream.}

\section{Branching and Merging}\label{branching-and-merging}

Branches are one of Git's most powerful features. A branch is an
independent line of development that diverges from the main project
history. You can experiment freely on a branch without affecting the
stable version of your code.

\textbf{Creating and switching to a new branch:}

\begin{Shaded}
\begin{Highlighting}[]
\FunctionTok{git}\NormalTok{ branch experiment}
\FunctionTok{git}\NormalTok{ checkout experiment}
\end{Highlighting}
\end{Shaded}

Or, in a single command:

\begin{Shaded}
\begin{Highlighting}[]
\FunctionTok{git}\NormalTok{ checkout }\AttributeTok{{-}b}\NormalTok{ experiment}
\end{Highlighting}
\end{Shaded}

\textbf{Working on the branch:}

Make your changes, stage, and commit as usual. The commits will only
exist on the \texttt{experiment} branch.

\begin{Shaded}
\begin{Highlighting}[]
\FunctionTok{git}\NormalTok{ add .}
\FunctionTok{git}\NormalTok{ commit }\AttributeTok{{-}m} \StringTok{"Test random forest approach"}
\end{Highlighting}
\end{Shaded}

\textbf{Merging back into main:}

When the experiment succeeds, merge it into the main branch:

\begin{Shaded}
\begin{Highlighting}[]
\FunctionTok{git}\NormalTok{ checkout main}
\FunctionTok{git}\NormalTok{ merge experiment}
\end{Highlighting}
\end{Shaded}

\textbf{Cleaning up:}

Once merged, delete the branch to keep the repository tidy:

\begin{Shaded}
\begin{Highlighting}[]
\FunctionTok{git}\NormalTok{ branch }\AttributeTok{{-}d}\NormalTok{ experiment}
\end{Highlighting}
\end{Shaded}

I found branching indispensable for data science work. When I wanted to
test whether a different variable selection strategy improved model
performance, I created a branch, ran the analysis, compared results, and
either merged or discarded the branch. The main branch always contained
the working version.

\section{Collaborating on GitHub}\label{collaborating-on-github}

One of the most common collaboration scenarios involves inviting a team
member to contribute to an existing repository. I encountered this when
I had been working solo on an analysis of data from the public
Alzheimer's Disease Neuroimaging Initiative (ADNI) repository, and the
project had grown complex enough that I needed help.

\textbf{Adding a collaborator (repository owner):}

\begin{enumerate}
\def\labelenumi{\arabic{enumi}.}
\tightlist
\item
  Log into GitHub and navigate to the repository (for example,
  \texttt{rgt47/x24}).
\item
  Select \textbf{Settings} in the top row of tabs.
\item
  Select \textbf{Collaborators} in the left panel.
\item
  Click the green \textbf{Add people} button.
\item
  Search for the collaborator's GitHub username (for example,
  \texttt{rgt4748}).
\item
  Select \textbf{Add rgt4748} to send the invitation.
\end{enumerate}

\textbf{Accepting the invitation (collaborator):}

\begin{enumerate}
\def\labelenumi{\arabic{enumi}.}
\tightlist
\item
  Log into GitHub.
\item
  Click the notification icon in the upper right corner.
\item
  Select the invitation (for example, ``Invitation to join rgt47/x24
  from rgt47'').
\item
  Click the green \textbf{Accept Invitation} button.
\end{enumerate}

\textbf{Cloning and contributing (collaborator):}

The collaborator clones the repository to their workstation:

\begin{Shaded}
\begin{Highlighting}[]
\FunctionTok{git}\NormalTok{ clone }\DataTypeTok{\textbackslash{}}
\NormalTok{  https://github.com/rgt47/x24.git}
\BuiltInTok{cd}\NormalTok{ x24}
\end{Highlighting}
\end{Shaded}

They create a branch for their work:

\begin{Shaded}
\begin{Highlighting}[]
\FunctionTok{git}\NormalTok{ checkout }\AttributeTok{{-}b}\NormalTok{ myedits}
\end{Highlighting}
\end{Shaded}

They make their changes (for example, editing \texttt{x24.Rmd}), then
stage, commit, and push:

\begin{Shaded}
\begin{Highlighting}[]
\FunctionTok{git}\NormalTok{ add .}
\FunctionTok{git}\NormalTok{ commit }\AttributeTok{{-}m} \StringTok{"Update analysis title"}
\FunctionTok{git}\NormalTok{ push origin myedits}
\end{Highlighting}
\end{Shaded}

\textbf{Creating a pull request:}

On GitHub, the collaborator clicks \textbf{Contribute}, then
\textbf{Open pull request}, enters a title and description, and submits
the request.

\textbf{Reviewing and merging (repository owner):}

The repository owner reviews the pull request, checks the \textbf{Files
changed} tab, and merges if the changes are acceptable. This workflow
keeps the main branch clean and provides a record of every contribution.

\section{Commit Message Best
Practices}\label{commit-message-best-practices}

I initially treated commit messages as an afterthought, writing things
like ``fixed stuff'' or ``updates.'' Over time, I learned that
well-crafted commit messages are one of the most valuable artifacts in a
project's history.

\subsection{Use the Imperative Mood}\label{use-the-imperative-mood}

A valuable practice involves writing commit messages as though the
commit, when applied, will perform a specific action. Construct your
message so that it logically completes the sentence: ``If applied, this
commit will\ldots{}''

Instead of:

\begin{Shaded}
\begin{Highlighting}[]
\FunctionTok{git}\NormalTok{ commit }\AttributeTok{{-}m} \StringTok{"Fixed the bug on the layout page"}
\end{Highlighting}
\end{Shaded}

Write:

\begin{Shaded}
\begin{Highlighting}[]
\FunctionTok{git}\NormalTok{ commit }\AttributeTok{{-}m} \StringTok{"Fix the bug on the layout page"}
\end{Highlighting}
\end{Shaded}

If this commit were applied, it would indeed fix the bug on the layout
page. The imperative mood keeps messages consistent and action-oriented.

\subsection{Explain ``What'' and ``Why,'' Not
``How''}\label{explain-what-and-why-not-how}

Limiting commit messages to ``what'' and ``why'' creates concise yet
informative explanations. Developers seeking the ``how'' can refer to
the code itself. Instead of describing implementation details, highlight
what was changed and the rationale behind it.

\subsection{The Angular Convention}\label{the-angular-convention}

The Angular project provides an exemplary model for structured commit
messages. Their convention uses specific prefixes that categorize each
commit:

\begin{itemize}
\tightlist
\item
  \texttt{feat:} -- a new feature
\item
  \texttt{fix:} -- a bug fix
\item
  \texttt{docs:} -- documentation changes
\item
  \texttt{style:} -- formatting, no code change
\item
  \texttt{refactor:} -- code restructuring
\item
  \texttt{test:} -- adding or updating tests
\item
  \texttt{chore:} -- maintenance tasks
\end{itemize}

By incorporating these prefixes, the commit history becomes a structured
resource for understanding the nature and purpose of each change. For
example:

\begin{Shaded}
\begin{Highlighting}[]
\FunctionTok{git}\NormalTok{ commit }\AttributeTok{{-}m} \StringTok{"feat: Add cross{-}validation to model"}
\FunctionTok{git}\NormalTok{ commit }\AttributeTok{{-}m} \StringTok{"fix: Correct off{-}by{-}one in data merge"}
\FunctionTok{git}\NormalTok{ commit }\AttributeTok{{-}m} \StringTok{"docs: Update README with new workflow"}
\end{Highlighting}
\end{Shaded}

I adopted a simplified version of this convention for my data science
projects and found that it made navigating project history substantially
easier.

\section{Viewing History and Recovering
Work}\label{viewing-history-and-recovering-work}

Git's history commands are essential for understanding how a project
evolved and for recovering from mistakes.

\textbf{Viewing the commit log:}

\begin{Shaded}
\begin{Highlighting}[]
\FunctionTok{git}\NormalTok{ log}
\FunctionTok{git}\NormalTok{ log }\AttributeTok{{-}{-}oneline} \AttributeTok{{-}{-}graph}
\end{Highlighting}
\end{Shaded}

The \texttt{-\/-oneline\ -\/-graph} flags produce a compact, visual
representation of the branch and merge history.

\textbf{Viewing a file from a specific commit:}

\begin{Shaded}
\begin{Highlighting}[]
\FunctionTok{git}\NormalTok{ show master:analysis.Rmd}
\end{Highlighting}
\end{Shaded}

This displays the file contents as they existed on the specified branch
or commit, without modifying your working directory.

\textbf{Restoring a file from another branch:}

\begin{Shaded}
\begin{Highlighting}[]
\FunctionTok{git}\NormalTok{ checkout main }\AttributeTok{{-}{-}}\NormalTok{ data/results.csv}
\end{Highlighting}
\end{Shaded}

\textbf{Recovering an earlier version of the entire project:}

\begin{Shaded}
\begin{Highlighting}[]
\FunctionTok{git}\NormalTok{ checkout abc1234}
\end{Highlighting}
\end{Shaded}

Where \texttt{abc1234} is the commit hash. This places you in a
``detached HEAD'' state where you can inspect the old version. To return
to the current state:

\begin{Shaded}
\begin{Highlighting}[]
\FunctionTok{git}\NormalTok{ checkout main}
\end{Highlighting}
\end{Shaded}

\textbf{Editing the Git configuration directly:}

For advanced users, the repository-level configuration file at
\texttt{.git/config} can be edited to adjust remote URLs, branch
tracking, and other settings.

\section{Troubleshooting Common
Issues}\label{troubleshooting-common-issues}

When working with Git, I encountered several recurring problems that
required specific solutions.

\textbf{Unrelated histories error:}

When merging repositories that were initialized independently, Git may
refuse with an ``unrelated histories'' error. The fix:

\begin{Shaded}
\begin{Highlighting}[]
\FunctionTok{git}\NormalTok{ pull }\AttributeTok{{-}{-}allow{-}unrelated{-}histories}
\end{Highlighting}
\end{Shaded}

\textbf{Large files accidentally committed:}

If you commit a large data file by mistake, it becomes part of the
repository history even after deletion. Use \texttt{.gitignore} to
prevent this:

\begin{Shaded}
\begin{Highlighting}[]
\ExtensionTok{*.csv}
\ExtensionTok{*.rds}
\ExtensionTok{data/raw/}
\end{Highlighting}
\end{Shaded}

\textbf{Merge conflicts:}

When two people edit the same lines of the same file, Git cannot
automatically merge the changes. Open the conflicted file, look for the
\texttt{\textless{}\textless{}\textless{}\textless{}\textless{}\textless{}\textless{}},
\texttt{=======}, and
\texttt{\textgreater{}\textgreater{}\textgreater{}\textgreater{}\textgreater{}\textgreater{}\textgreater{}}
markers, resolve the conflict manually, then stage and commit.

\textbf{Authentication failures:}

If \texttt{git\ push} fails with a permission error, verify that your
SSH key is added to the agent:

\begin{Shaded}
\begin{Highlighting}[]
\BuiltInTok{eval} \StringTok{"}\VariableTok{$(}\FunctionTok{ssh{-}agent} \AttributeTok{{-}s}\VariableTok{)}\StringTok{"}
\FunctionTok{ssh{-}add}\NormalTok{ \textasciitilde{}/.ssh/id\_ed25519}
\end{Highlighting}
\end{Shaded}

\section{Verification}\label{verification}

After installing and configuring Git, confirm that the setup is working
correctly.

\begin{Shaded}
\begin{Highlighting}[]
\FunctionTok{git} \AttributeTok{{-}{-}version}
\end{Highlighting}
\end{Shaded}

\begin{Shaded}
\begin{Highlighting}[]
\FunctionTok{git}\NormalTok{ config user.name }\KeywordTok{\&\&} \FunctionTok{git}\NormalTok{ config user.email}
\end{Highlighting}
\end{Shaded}

\begin{Shaded}
\begin{Highlighting}[]
\FunctionTok{ssh} \AttributeTok{{-}T}\NormalTok{ git@github.com}
\end{Highlighting}
\end{Shaded}

\begin{Shaded}
\begin{Highlighting}[]
\FunctionTok{mkdir}\NormalTok{ /tmp/test{-}repo }\KeywordTok{\&\&} \BuiltInTok{cd}\NormalTok{ /tmp/test{-}repo }\KeywordTok{\&\&} \DataTypeTok{\textbackslash{}}
  \FunctionTok{git}\NormalTok{ init }\KeywordTok{\&\&} \FunctionTok{touch}\NormalTok{ README.md }\KeywordTok{\&\&} \DataTypeTok{\textbackslash{}}
  \FunctionTok{git}\NormalTok{ add . }\KeywordTok{\&\&} \FunctionTok{git}\NormalTok{ commit }\AttributeTok{{-}m} \StringTok{"test"} \KeywordTok{\&\&} \DataTypeTok{\textbackslash{}}
  \FunctionTok{git}\NormalTok{ log }\AttributeTok{{-}{-}oneline}
\end{Highlighting}
\end{Shaded}

If \texttt{git\ -\/-version} returns 2.30+, the config shows your name
and email, SSH authenticates successfully, and the test commit appears
in the log, the setup is complete.

\section{Daily Workflow}\label{daily-workflow}

{\def\LTcaptype{none} % do not increment counter
\begin{longtable}[]{@{}ll@{}}
\toprule\noalign{}
Task & Command \\
\midrule\noalign{}
\endhead
\bottomrule\noalign{}
\endlastfoot
Check status & \texttt{git\ status} \\
Stage files & \texttt{git\ add\ analysis.R\ data/clean.csv} \\
Commit changes & \texttt{git\ commit\ -m\ "feat:\ Add\ model"} \\
Push to remote & \texttt{git\ push\ origin\ main} \\
Pull latest & \texttt{git\ pull\ origin\ main} \\
Create branch & \texttt{git\ checkout\ -b\ experiment} \\
Merge branch &
\texttt{git\ checkout\ main\ \&\&\ git\ merge\ experiment} \\
\end{longtable}
}

\subsection{Things to Watch Out For}\label{things-to-watch-out-for}

\begin{enumerate}
\def\labelenumi{\arabic{enumi}.}
\item
  \textbf{Never commit credentials or API keys.} Add \texttt{.env},
  \texttt{.Renviron}, and credential files to \texttt{.gitignore} before
  your first commit. Once a secret is in Git history, removing it
  completely requires rewriting history.
\item
  \textbf{Avoid committing large binary files.} CSV files over 50 MB,
  RDS objects, and image datasets do not belong in Git. Use Git LFS or
  external storage for large data.
\item
  \textbf{Pull before you push.} When collaborating, always
  \texttt{git\ pull} before starting new work. Pushing without pulling
  first leads to merge conflicts that could have been avoided.
\item
  \textbf{Do not force-push to shared branches.} Running
  \texttt{git\ push\ -\/-force} on \texttt{main} rewrites history for
  all collaborators. Reserve force-push for personal branches only.
\item
  \textbf{Commit frequently, push regularly.} Small, focused commits are
  easier to review, easier to revert, and create a more useful history
  than large, infrequent commits.
\end{enumerate}

\section{Uninstall / Rollback}\label{uninstall-rollback}

To remove Git configuration without uninstalling Git itself:

\begin{Shaded}
\begin{Highlighting}[]
\FunctionTok{git}\NormalTok{ config }\AttributeTok{{-}{-}global} \AttributeTok{{-}{-}unset}\NormalTok{ user.name}
\FunctionTok{git}\NormalTok{ config }\AttributeTok{{-}{-}global} \AttributeTok{{-}{-}unset}\NormalTok{ user.email}
\FunctionTok{git}\NormalTok{ config }\AttributeTok{{-}{-}global} \AttributeTok{{-}{-}unset}\NormalTok{ init.defaultBranch}
\FunctionTok{git}\NormalTok{ config }\AttributeTok{{-}{-}global} \AttributeTok{{-}{-}unset}\NormalTok{ core.editor}
\end{Highlighting}
\end{Shaded}

To remove SSH keys used for GitHub:

\begin{Shaded}
\begin{Highlighting}[]
\FunctionTok{rm}\NormalTok{ \textasciitilde{}/.ssh/id\_ed25519 \textasciitilde{}/.ssh/id\_ed25519.pub}
\end{Highlighting}
\end{Shaded}

To remove Git entirely on macOS, uninstall the Command Line Tools. On
Ubuntu:

\begin{Shaded}
\begin{Highlighting}[]
\FunctionTok{sudo}\NormalTok{ apt remove git}
\end{Highlighting}
\end{Shaded}

To remove version control from a single project, delete its
\texttt{.git} directory:

\begin{Shaded}
\begin{Highlighting}[]
\FunctionTok{rm} \AttributeTok{{-}rf}\NormalTok{ .git}
\end{Highlighting}
\end{Shaded}

This is irreversible and destroys all project history.

\begin{figure}[H]

{\centering \pandocbounded{\includegraphics[keepaspectratio]{media/images/ambiance3.jpg}}

}

\caption{UCSD Geisel Library, a landmark of research and learning}

\end{figure}%

\emph{Disciplined version control practices support the same rigor that
characterizes scholarly research.}

\section{What Did We Learn?}\label{what-did-we-learn}

\subsection{Lessons Learnt}\label{lessons-learnt}

\textbf{Conceptual Understanding:}

\begin{itemize}
\tightlist
\item
  Version control is not merely a backup system; it is a structured
  record of a project's intellectual history, capturing the reasoning
  behind every change.
\item
  Branching transforms experimentation from a risky activity into a
  safe, reversible process, which is particularly valuable in
  exploratory data analysis.
\item
  The distributed nature of Git means that every collaborator holds a
  complete copy of the project, eliminating single points of failure.
\item
  Commit messages, when written with care, become a form of
  documentation that is often more useful than formal project notes.
\end{itemize}

\textbf{Technical Skills:}

\begin{itemize}
\tightlist
\item
  The core workflow (\texttt{add}, \texttt{commit}, \texttt{push},
  \texttt{pull}) covers the vast majority of daily version control needs
  for a solo analyst.
\item
  SSH key authentication eliminates password friction and is more secure
  than HTTPS token-based access for regular use.
\item
  The \texttt{git\ log\ -\/-oneline\ -\/-graph} command provides a rapid
  overview of project history and branching structure.
\item
  The \texttt{.gitignore} file is the first line of defense against
  accidentally committing sensitive or oversized files.
\end{itemize}

\textbf{Gotchas and Pitfalls:}

\begin{itemize}
\tightlist
\item
  Forgetting to add a \texttt{.gitignore} before the first commit can
  result in large data files entering the repository history
  permanently.
\item
  Using \texttt{git\ add\ .} without reviewing staged files is a common
  source of accidental credential leaks.
\item
  Merge conflicts in data files (CSV, JSON) are particularly difficult
  to resolve because the conflict markers corrupt the file format.
\item
  Working directly on \texttt{main} instead of a branch removes the
  safety net that makes Git valuable for experimentation.
\end{itemize}

\subsection{Limitations}\label{limitations}

\begin{itemize}
\tightlist
\item
  This guide covers only the command-line interface. GUI tools like
  GitKraken, Sourcetree, and the RStudio Git pane provide alternative
  workflows that some users may prefer.
\item
  The collaboration scenario assumes a small team with direct repository
  access. Larger projects typically use fork-based workflows, which are
  not covered here.
\item
  Git tracks text files effectively but is poorly suited for binary
  files such as images, compiled documents, and large datasets.
\item
  This tutorial does not cover Git internals (the object model,
  packfiles, or the reflog), which are important for advanced
  troubleshooting.
\item
  The commit message conventions described here reflect one widely
  adopted standard (Angular). Other projects use different conventions,
  and consistency within a team matters more than which specific
  convention is chosen.
\item
  Windows-specific configuration (line endings, path length limits,
  credential managers) is not addressed.
\end{itemize}

\subsection{Opportunities for
Improvement}\label{opportunities-for-improvement}

\begin{enumerate}
\def\labelenumi{\arabic{enumi}.}
\tightlist
\item
  Explore Git hooks (pre-commit, pre-push) to automate code linting,
  testing, and credential scanning before changes enter the repository.
\item
  Investigate Git LFS (Large File Storage) for managing datasets that
  exceed Git's practical size limits.
\item
  Adopt a branching strategy such as Git Flow or trunk-based development
  for projects with multiple active contributors.
\item
  Integrate Git with continuous integration services (GitHub Actions,
  GitLab CI) to automatically run tests and render reports on every
  push.
\item
  Learn interactive rebase (\texttt{git\ rebase\ -i}) for cleaning up
  commit history before merging feature branches.
\item
  Set up signed commits using GPG keys to verify authorship in
  security-sensitive research contexts.
\end{enumerate}

\section{Wrapping Up}\label{wrapping-up}

Setting up Git for data science work requires a modest initial
investment---perhaps an afternoon of configuration and practice---but
the returns are substantial. Every analysis becomes recoverable, every
experiment becomes safe to try, and every collaboration becomes
structured.

What I learned most clearly from adopting Git is that version control
changes the way you think about your work. Once every change is tracked,
you become more willing to experiment, more disciplined about
documentation, and more confident that nothing important will be lost.

For readers getting started with Git, my advice is straightforward:

\begin{itemize}
\tightlist
\item
  Start with the core workflow: \texttt{init}, \texttt{add},
  \texttt{commit}, \texttt{push}.
\item
  Use branches for every experiment, no matter how small.
\item
  Write commit messages in the imperative mood with a clear prefix.
\item
  Add a \texttt{.gitignore} before your first commit.
\item
  Pull before you push when collaborating.
\end{itemize}

The commit history you build today is the documentation your future self
will rely on. Make it a habit to create messages that stand as
informative, concise, and consistent narratives of your project's
evolution.

\section{See Also}\label{see-also}

\textbf{Related blog posts:}

\begin{itemize}
\tightlist
\item
  \href{../24-setupdotfilesongithub/setupdotfilesongithub/analysis/report/index.qmd}{Setting
  Up Dotfiles on GitHub} -- Managing configuration files with Git
\item
  \href{../01-configtermzsh/configtermzsh/analysis/report/index.qmd}{Configuring
  the Terminal for Data Science} -- Terminal setup that complements a
  Git workflow
\end{itemize}

\textbf{Key external resources:}

\begin{itemize}
\tightlist
\item
  \href{https://git-scm.com/book/en/v2}{Pro Git Book} -- The definitive
  free reference by Scott Chacon and Ben Straub
\item
  \href{https://git-scm.com/doc}{Git Official Documentation} -- Command
  reference and manual pages
\item
  \href{https://skills.github.com/}{GitHub Skills} -- Interactive
  courses for learning GitHub workflows
\item
  \href{https://www.atlassian.com/git/tutorials}{Atlassian Git
  Tutorials} -- Clear, well-structured guides for all skill levels
\item
  \href{https://learngitbranching.js.org/}{Learn Git Branching} --
  Interactive visualization tool for understanding branches
\item
  \href{https://www.git-tower.com/learn/}{Git Tower Learning Resources}
  -- Tutorials and cheat sheets
\item
  \href{https://www.gitkraken.com/git-cheat-sheet}{GitKraken Git Cheat
  Sheet} -- Quick command reference
\item
  Bryan, J. (2018). Excuse Me, Do You Have a Moment to Talk About
  Version Control? \emph{The American Statistician}, 72(1), 20--27.
  \href{https://doi.org/10.1080/00031305.2017.1399928}{DOI:
  10.1080/00031305.2017.1399928}
\item
  Bonfim, G. (2023).
  \href{https://gabrielbonfim1.medium.com/git-basics-all-you-need-to-know-as-a-new-developer-f4d905e4e144}{Git
  Basics---All You Need to Know as a New Developer}. Medium.
\end{itemize}

\section{Reproducibility}\label{reproducibility}

This post does not involve computational analysis, so there is no
analysis pipeline to reproduce. The Git commands shown throughout the
post can be executed in any standard terminal environment with Git
installed.

\textbf{Environment used for this post:}

\begin{Shaded}
\begin{Highlighting}[]
\FunctionTok{git} \AttributeTok{{-}{-}version}
\end{Highlighting}
\end{Shaded}

\textbf{Minimum requirements:}

\begin{itemize}
\tightlist
\item
  Git 2.30 or later
\item
  macOS 12+ or Ubuntu 20.04+
\item
  A GitHub account with SSH key configured
\end{itemize}

All commands were tested on macOS with Zsh and Git 2.43.0.

\section{Let's Connect}\label{lets-connect}

\begin{itemize}
\tightlist
\item
  \textbf{GitHub:} \href{https://github.com/rgt47}{rgt47}
\item
  \textbf{Twitter/X:} \href{https://twitter.com/rgt47}{@rgt47}
\item
  \textbf{LinkedIn:} \href{https://linkedin.com/in/rgthomaslab}{Ronald
  Glenn Thomas}
\item
  \textbf{Email:} \href{https://rgtlab.org/contact}{rgtlab.org/contact}
\end{itemize}

I would enjoy hearing from you if:

\begin{itemize}
\tightlist
\item
  You spot an error or a better approach to any of the code in this
  post.
\item
  You have suggestions for topics you would like to see covered.
\item
  You want to discuss R programming, data science, or reproducible
  research.
\item
  You have questions about anything in this tutorial.
\item
  You just want to say hello and connect.
\end{itemize}




\end{document}
